%% Compile: pdflatex -interaction=nonstopmode methods-note.tex  (twice, for the references)
\documentclass[11pt]{article}
\usepackage[utf8]{inputenc}
\usepackage[T1]{fontenc}
\usepackage{amsmath,amssymb,amsthm}
\usepackage{booktabs}
\usepackage[margin=2.4cm]{geometry}
\usepackage[colorlinks=true,linkcolor=blue,citecolor=blue]{hyperref}
\newtheorem{proposition}{Proposition}
\newtheorem{remark}{Remark}
\DeclareMathOperator{\Var}{Var}
\DeclareMathOperator{\Cov}{Cov}
\DeclareMathOperator{\Corr}{Corr}
\newcommand{\bbeta}{\boldsymbol\beta}
\newcommand{\bSigma}{\boldsymbol\Sigma}
\newcommand{\bmu}{\boldsymbol\mu}
\newcommand{\bff}{\mathbf f}
\newcommand{\bR}{\mathbf R}
\newcommand{\br}{\mathbf r}

\title{A specification layer for factor-based market simulation\\
\large Derivations for the CVineMarketGen notebooks}
\author{Mamadou Yamar Thioub}
\date{September 2026}

\begin{document}
\maketitle

\begin{abstract}
\noindent An analyst who wants scenarios for a universe of assets rarely has a
complete description of it. They have a history for some assets, a mean or a
volatility for others, exposures for an asset class, a correlation they insist
on, a view on a driver. This note states the rule that turns any such partial
description into one simulator, and derives the boundaries of what can be
asked: the feasible residual variance, the feasible residual shape, the range
of a pair correlation, the largest correlation an asset can have with a factor,
the completion of the correlations of a factor with no history. It closes with
two selection questions: which copula family a pair of near-orthogonal factors
can identify, and how the dynamics of a factor are chosen. Numbers appear only
where one makes a statement concrete; the notebooks carry the numbers.
\end{abstract}

\section{The object}

The simulator draws $K$ common factors and maps each simulated factor month to
$N$ assets,
\begin{equation}\label{eq:model}
  r_j = \alpha_j + \bbeta_j' \bff + e_j, \qquad j = 1, \dots, N,
\end{equation}
with $\bff$ the factor vector, $e_j$ independent of $\bff$ and, except where a
pair is targeted, of the other residuals. Writing $\mathbf B$ for the $N \times
K$ matrix of betas and $\mathbf D$ for the diagonal matrix of residual
variances, the implied covariance is
\begin{equation}\label{eq:cov}
  \bSigma_r = \mathbf B\, \bSigma_f\, \mathbf B' + \mathbf D ,
\end{equation}
which is what makes the exercise possible at all: $K(K+1)/2$ factor
correlations and $NK$ betas in place of $N(N-1)/2$ asset correlations, and an
asset with no history reduced to a row of $\mathbf B$ and one variance.

The factor vector is drawn by a generator that reproduces the factors' means,
volatilities, skewnesses, kurtoses, correlations, and, in the C-vine case,
their pairwise tail dependence. The specification layer sits above it: it
decides what the generator is asked for, and what each asset carries.

\section{The four dimensions of knowledge, and one order}

What an analyst can supply falls in four groups: a \emph{history} (of an asset,
of a factor, of neither); \emph{moments} (a mean, a volatility, a skewness, a
kurtosis, for an asset or for a factor); \emph{exposures} (a beta, or a
correlation with a factor, stated for an asset or for a whole class); and
\emph{pair correlations} between two assets or two factors. Any subset may be
given, in any combination.

The rule is a fixed order, each step taking the previous ones as given:
\begin{enumerate}
  \item \textbf{betas} --- stated values; correlations with a factor solved for
        the beta; regression on the factors marked as fitted; zero elsewhere;
  \item \textbf{volatility} --- the residual variance is the target variance
        minus the systematic variance, or the history's residual variance when
        no target is given;
  \item \textbf{shape} --- the residual cumulants are the target cumulants
        minus the systematic cumulants, or the history's residual shape;
  \item \textbf{pair correlations} --- the residual correlation that delivers
        the target, inside its feasible range;
  \item \textbf{mean} --- the intercept absorbs it.
\end{enumerate}
The order is forced: the residual's scale cannot be computed before the betas,
its shape not before its scale, a pair's room not before both, and the
intercept is free of all of them. An empty cell is not a missing value but an
instruction: estimate it. What is estimated is the model's own quantity --- the
beta, the residual, the intercept --- and not the asset's moment, which is why a
view on a factor reaches an asset whose cells are empty.

\section{Betas}

\subsection{A correlation with a factor as a linear constraint}

Let $S$ be the set of factors for which a correlation $\rho_k$ with asset $j$
is stated rather than a beta. Since $\Cov(r_j, f_k) = (\bSigma_f
\bbeta_j)_k$, the constraint $\Corr(r_j, f_k) = \rho_k$ for $k \in S$ reads
\begin{equation}\label{eq:corr_beta}
  \bSigma_{SS}\, \bbeta_S \;=\; \sigma_j\, \big(\rho_k \sigma_{f_k}\big)_{k \in S}
     \;-\; \bSigma_{S,-S}\, \bbeta_{-S} ,
\end{equation}
a linear system in $\bbeta_S$ once $\sigma_j$ and the other betas are known.
When the asset has a volatility of its own, $\sigma_j$ is given and
\eqref{eq:corr_beta} is solved once. When it does not, $\sigma_j$ is itself a
function of the betas, and \eqref{eq:corr_beta} is iterated with
\begin{equation}\label{eq:fixed_point}
  \sigma_j^2 \;=\; \bbeta_j' \bSigma_f \bbeta_j + \Var(e_j), \qquad
  e_j \;=\; r_j - \bbeta_j' \bff ,
\end{equation}
the residual being re-estimated at each pass. The iteration converges quickly
in practice, but not every target is reachable.

\subsection{The saturation bound}

\begin{proposition}\label{prop:sat}
Let asset $j$ have a history, no volatility target, and one factor $f$ with
$\Corr(r_j, f) = \rho_0$ in the sample. If the beta on $f$ is set to $b$ and
the residual is the empirical residual $e = r_j - bf$, then the correlation the
model delivers is maximised at $b^\star = \sigma_{r_j} / (\rho_0 \sigma_f)$ and
equals
\begin{equation}\label{eq:rho_max}
  \rho_{\max} \;=\; \frac{1}{\sqrt{2 - R_j^2}}, \qquad R_j^2 = \rho_0^2 .
\end{equation}
\end{proposition}

\begin{proof}
Write $x = b \sigma_f / \sigma_{r_j}$. The residual variance is
$\Var(e) = \sigma_{r_j}^2 - 2b\,\Cov(r_j, f) + b^2 \sigma_f^2$, so by
\eqref{eq:fixed_point} the asset's variance under the imposed beta is
$\sigma_j^2 = 2b^2\sigma_f^2 - 2b\,\Cov(r_j,f) + \sigma_{r_j}^2$, and
\[
  \rho(x) \;=\; \frac{b \sigma_f^2}{\sigma_f \sigma_j}
           \;=\; \frac{x}{\sqrt{1 - 2\rho_0 x + 2x^2}} .
\]
Setting the derivative of $\rho^2$ to zero gives $2x(1 - \rho_0 x) = 0$, hence
$x^\star = 1/\rho_0$ and $\rho(x^\star) = (1/\rho_0) / \sqrt{2/\rho_0^2 - 1} =
1/\sqrt{2 - \rho_0^2}$.
\end{proof}

\begin{remark}
The bound is the price of leaving the volatility cell empty: raising the beta
raises the covariance, but it raises the asset's own variance twice as fast,
because the residual no longer fits. With several factors, $R_j^2$ is the share
of variance the factors explain and \eqref{eq:rho_max} is an upper bound rather
than an identity. A volatility in the assets sheet pins $\Var(e_j)$ through
$\sigma_{e}^2 = v^2 - \bbeta'\bSigma_f\bbeta$ and removes the bound: any
correlation up to the variance-feasibility limit is then attainable.
\end{remark}

\section{Volatility}

With a target volatility $v$ and the betas fixed,
\begin{equation}\label{eq:vol}
  \sigma_e^2 \;=\; v^2 - \bbeta'\bSigma_f\bbeta \;>\; 0 ,
\end{equation}
so a target is feasible if and only if it exceeds the volatility the betas
alone produce. This is the one refusal an analyst meets most often, and it is
informative: it says the exposures and the volatility are inconsistent, and the
refusal names the two largest contributors to $\bbeta'\bSigma_f\bbeta$ so that
one knows which beta to lower. With no target, $\sigma_e$ is the sample residual
standard deviation and the asset's volatility follows from \eqref{eq:cov}: a
view on a factor then moves the asset's volatility, as it should.

\section{Shape}

\subsection{The cumulant rule}

Cumulants of independent random variables add. With $s$ and $k$ the target
skewness and kurtosis of $r_j$, $v$ its volatility and $\kappa_3, \kappa_4$ the
third and fourth cumulants,
\begin{equation}\label{eq:cumulants}
  \kappa_{3,e} \;=\; s\,v^3 - \kappa_3(\bbeta'\bff), \qquad
  \kappa_{4,e} \;=\; (k - 3)\,v^4 - \kappa_4(\bbeta'\bff),
\end{equation}
and the residual's standardised shape is $s_e = \kappa_{3,e}/\sigma_e^3$,
$k_e = 3 + \kappa_{4,e}/\sigma_e^4$. The systematic cumulants are read off a
large pilot draw of the fitted factor generator, which costs nothing and avoids
assuming anything about the factors' joint law.

Two things follow. First, the residual must carry the whole difference between
what is asked and what the factors already give, so the smaller the residual
share $\sigma_e^2/v^2$, the larger the shape the residual must have: dividing
by $\sigma_e^3$ and $\sigma_e^4$ magnifies the demand. Second, the demand may
leave the family of distributions the residual is drawn from.

\subsection{The Johnson $S_U$ feasible region}

The residual is drawn from a Johnson $S_U$ law,
$X = \xi + \lambda \sinh\!\big((Z - \gamma)/\delta\big)$ with $Z$ standard
normal, whose $(s, k)$ region is bounded below by the lognormal line: with
$w = e^{1/\delta^2}$,
\begin{equation}\label{eq:jsu_boundary}
  |s| \;=\; (w + 2)\sqrt{w - 1}, \qquad
  k \;=\; w^4 + 2w^3 + 3w^2 - 3 .
\end{equation}
Given $|s|$, the first equation has one root $w > 1$ (the left side is
increasing) and the second gives the smallest kurtosis the family reaches at
that skewness. The floor is $3$ at $s = 0$ and grows quickly: about $3.8$ at
$|s| = 0.6$ and $7.6$ at $|s| = 1.5$. A pair $(s_e, k_e)$ below the line is
refused, and the refusal reports $s_e$, $k_e$, the residual share and the shape
the factors already give, which together say whether to soften the target or to
lower the exposures.

\begin{remark}
The same floor acts silently in the other direction. When no shape is given and
the sample residual of an asset is flatter than the family allows --- a residual
kurtosis of $2.9$ at a skewness of $-0.3$, say --- the kurtosis is raised to the
floor and the asset is flagged. The change is small, and hiding it would be
worse than making it.
\end{remark}

\subsection{How precisely a shape is delivered}

A mean and a volatility are matched by construction, up to Monte Carlo error. A
shape is not: it is matched through a fitted marginal whose parameters solve
\eqref{eq:cumulants} at the third and fourth moment, and the sample skewness
and kurtosis of a finite draw of a heavy-tailed variable are themselves noisy.
In the notebooks, a target skewness of $-0.5$ and kurtosis of $5$ on an asset
whose residual is two thirds of its variance lands within about one tenth in
skewness and half a point in kurtosis across seeds. A shape is a target, not an
identity, and should be read at that resolution.

\section{The correlation between two assets}

\subsection{The eligible range}

From \eqref{eq:model}, for two assets $i \ne j$,
\begin{equation}\label{eq:pair}
  \rho_{ij} \;=\; \underbrace{\frac{\bbeta_i' \bSigma_f \bbeta_j}{\sigma_i \sigma_j}}_{\rho^{\text{sys}}_{ij}}
   \;+\; \rho_{e,ij}\sqrt{(1 - R_i^2)(1 - R_j^2)} ,
\end{equation}
where $R_i^2 = \bbeta_i'\bSigma_f\bbeta_i / \sigma_i^2$. The betas are already
fixed by step 1, so the only free quantity is $\rho_{e,ij} \in [-1, 1]$, and
\begin{equation}\label{eq:range}
  \rho_{ij} \in \Big[\rho^{\text{sys}}_{ij} - w_{ij},\;
                     \rho^{\text{sys}}_{ij} + w_{ij}\Big], \qquad
  w_{ij} = \sqrt{(1 - R_i^2)(1 - R_j^2)} ,
\end{equation}
with the residual correlation that delivers a feasible target given by
\begin{equation}\label{eq:rho_e}
  \rho_{e,ij} \;=\; \frac{\rho_{ij}\,\sigma_i\sigma_j - \bbeta_i'\bSigma_f\bbeta_j}
                         {\sigma_{e,i}\,\sigma_{e,j}} .
\end{equation}
Two well-explained assets have almost no room; two assets the factors barely
touch have almost the whole interval. This is the honest answer to ``can I set
the correlation between these two funds'': partly, and how much is a number you
can compute before asking.

\subsection{Several pairs at once}

Each target writes one off-diagonal entry of the residual correlation matrix
$\bR_e$ through \eqref{eq:rho_e}. Individually feasible targets may be jointly
infeasible, and the condition is the one the matrix imposes: $\bR_e \succeq 0$.
The smallest eigenvalue is checked once every pair has been written, and a
negative one is reported with the assets involved rather than projected away,
because a projection would silently move targets the analyst stated.

\section{A factor with no history}

A factor known only by a view needs a mean and a volatility --- there is nothing
to estimate them from --- and its skewness and kurtosis fall back to the normal.
Its correlations are the real question: an analyst will state one or two, and
$K$ are needed.

\begin{proposition}\label{prop:completion}
Let $\bR$ be the correlation matrix of the $K$ existing factors, $S$ the set of
factors whose correlation $\br_S$ with the new factor is stated, and $U$ the
rest. Among all $\br_U$ making the extended matrix positive semidefinite, the
determinant of that matrix is maximised by
\begin{equation}\label{eq:completion}
  \br_U \;=\; \bR_{US}\, \bR_{SS}^{-1}\, \br_S ,
\end{equation}
which is also the unique completion whose partial correlations between the new
factor and $U$ given $S$ are zero. A completion exists if and only if
$\br_S' \bR_{SS}^{-1} \br_S \le 1$.
\end{proposition}

\begin{proof}
The extended matrix has determinant $\det(\bR)\,(1 - \br'\bR^{-1}\br)$ with
$\br = (\br_S, \br_U)$, so maximising it is minimising the quadratic form
$\br'\bR^{-1}\br$ over $\br_U$. Writing $\mathbf Q = \bR^{-1}$ in blocks, the
first-order condition is $\mathbf Q_{US}\br_S + \mathbf Q_{UU}\br_U = 0$, hence
$\br_U = -\mathbf Q_{UU}^{-1}\mathbf Q_{US}\br_S$; the block-inverse identities
$\mathbf Q_{UU} = (\bR_{UU} - \bR_{US}\bR_{SS}^{-1}\bR_{SU})^{-1}$ and
$\mathbf Q_{US} = -\mathbf Q_{UU}\bR_{US}\bR_{SS}^{-1}$ turn this into
\eqref{eq:completion}. The conditional covariance of the new factor with $U$
given $S$ is $\br_U - \bR_{US}\bR_{SS}^{-1}\br_S$, zero exactly at
\eqref{eq:completion}. At the optimum $\br'\bR^{-1}\br = \br_S'\bR_{SS}^{-1}\br_S$,
so the extended matrix is positive semidefinite for some $\br_U$ if and only if
that quantity does not exceed one.
\end{proof}

The completion adds nothing to what was said: the new factor is related to the
factors nobody mentioned only through the ones that were. The feasibility
condition is also the whole story behind the refusal an analyst meets here. If
two existing factors correlate $0.61$ with each other, a new factor cannot
correlate $+0.9$ with one and $-0.9$ with the other: the quadratic form is then
$4.6$, far above one.

Once the correlations are settled, the new factor is attached to the fitted
generator by a Gaussian conditional draw. Let $\mathbf z$ be the normal scores
of the simulated existing factors and $\br$ the vector of target correlations.
The new factor's normal score is drawn as
\begin{equation}\label{eq:attach}
  z_{new} \mid \mathbf z \;\sim\; \mathcal N\!\big(\br'\bR^{-1}\mathbf z,\;
     1 - \br'\bR^{-1}\br\big),
\end{equation}
and mapped through the quantile function of its own marginal. The existing
factors keep the joint law the vine gave them --- the attachment does not touch
them --- and the new factor is joined to them by a Gaussian dependence, which is
the least one can assume when only correlations were stated.

\section{Why the copula layer is quiet on near-orthogonal factors}

A factor model is built so that its factors are nearly uncorrelated, and a
copula is then hard to identify. The reason is a matter of orders. Expand a
bivariate copula density in the Hermite basis of the normal scores; for the
Gaussian copula, Mehler's formula gives
\begin{equation}\label{eq:mehler}
  c_\rho(u,v) \;=\; \sum_{m \ge 0} \frac{\rho^m}{m!}\, H_m(x)\,H_m(y),
  \qquad x = \Phi^{-1}(u),\; y = \Phi^{-1}(v).
\end{equation}
Every one-parameter family reduces to the independence copula at zero
dependence, and matching two families on the Pearson correlation $\rho$ forces
their first Hermite coefficient to agree. Families therefore differ only from
the second coefficient on: their densities differ by $O(\rho^2)$, and the
expected log-likelihood gap per observation, a Kullback-Leibler divergence
between two densities at distance $O(\rho^2)$, is $O(\rho^4)$.

Put numbers on it. At $\rho = 0.1$ and $n = 230$ monthly observations, the
expected gap is of order $n\rho^4 \approx 0.02$ in log-likelihood, while the BIC
penalty for one extra parameter is $\tfrac12\log n \approx 2.7$ and the sampling
standard deviation of the gap is larger still. No selection rule can separate
the families there, and none should pretend to. The consequence for practice is
not that the copula is useless but that it earns its place elsewhere: between
asset classes, where correlations are $0.5$ and above and a Student t or a
rotated Gumbel is chosen against a Gaussian with a clear margin, and in the
tails of a regime model. At the premia level, a Gaussian dependence with the
right marginals is an honest answer.

Whether a pair's dependence is asymmetric is a separate question, and the
package answers it with a statistic rather than a family fit: the mean
exceedance correlation below $-1$ standard deviation minus the mean above $+1$,
computed on normal scores, with a bootstrap interval. A pair whose interval
contains zero is given a symmetric family and nothing else is read into it.

\section{Choosing the dynamics of a factor}

A one-period generator has no memory, and a factor's history has one. The
dynamics of a factor is chosen among a grid of candidates --- a constant or an
AR(1) mean, a constant variance, a GARCH, a GJR-GARCH or an EGARCH at orders up
to two, a Gaussian hidden Markov model with two or three regimes --- by one rule
in two steps.

\paragraph{Step 1: goodness of fit.} Every candidate is filtered, and its
generalized error is mapped to a uniform,
\begin{equation}\label{eq:rosenblatt}
  v_t \;=\; F\!\left(\varepsilon_t \mid \mathcal F_{t-1}; \hat{\boldsymbol\theta}\right) ,
\end{equation}
the Rosenblatt transform of the observation under the fitted model. If the
model is correctly specified, $(v_t)$ is i.i.d.\ uniform on $(0,1)$. The
discrepancy is measured by the Cramér-von Mises statistic
\begin{equation}\label{eq:cvm}
  S_n \;=\; \frac{1}{12n} + \sum_{i=1}^{n}\left(v_{(i)} - \frac{i - 1/2}{n}\right)^{\!2} ,
\end{equation}
whose null distribution depends on the estimated parameters, so the $p$-value
is obtained by parametric bootstrap: $B$ series are simulated from the fitted
model, each is refitted, its own statistic computed, and the $p$-value is the
share of bootstrap statistics above the observed one. A candidate is admissible
when that $p$-value is at least the level. For a GARCH filter whose innovation
is a Johnson $S_U$ --- what the generator actually simulates --- the uniforms are
taken through that marginal, and the bootstrap refits the marginal as well;
testing such a model against a Gaussian $F$ would reject it for a defect it
does not have.

\paragraph{Step 2: BIC among the admissible.} Among the candidates that pass,
the one with the lowest BIC is kept. The order matters and is not
interchangeable: the test asks whether a model is adequate, an information
criterion only ranks models against each other and will happily return the best
of a set of bad ones. Ljung-Box statistics on the generalized error and on its
square, and an ARCH-LM statistic, are reported alongside, as diagnostics of
\emph{what} a rejected model got wrong. They do not enter the rule.

If no candidate passes, the lowest BIC overall is returned with a warning
saying so, which is the honest output: the grid does not contain a model this
series cannot reject.

\section{What the rule does not do}

Three limits are worth stating plainly. The betas are estimated by ordinary
least squares on the factors marked as fitted, with no shrinkage and no
variable selection: an asset with a short history and many factors will
overfit, and the remedy in the specification layer is to state fewer factors
per class rather than to regress more. The residuals are independent except
where a pair is targeted, so a cluster of assets that move together for a
reason no factor captures needs one pair per link, and the joint
positive-semidefinite condition of \S6.2 will eventually bind. And the
scenarios are exchangeable months: the layer above, a filter per factor with
its errors drawn jointly, is what gives them an order in time.

\end{document}
